\scanpage{280}
\mspage{230}
On le voit sur le diagramme plus complet
\diagram{\begin{tikzcd}[ampersand replacement=\&,column sep=large,row sep=large]
\& X\times I \arrow[r,"f'"] \arrow[dd,dashed,bend left=48,"{\mathrm{pr}_1}" description] \& Z(f)\arrow[dd,bend left=40,dashed,"p_f"]\\
X\arrow[ur,dashed,"{\mathrm{inj}_1^X}"]\arrow[r,hook,"{\unsure{\mathrm{inj}|i}_1^X}"]
\& X\amalg X \arrow[u,hook,"i_X"'] \arrow[r,"{f\amalg\mathrm{id}_X}"]
\arrow[ur,phantom,"{\substack{\text{cocart et}\\\text{cart}}}" description]
\& Y\amalg X \arrow[u,hook,"j"']\\
\& X \arrow[u,hook,"{\unsure{i|\mathrm{inj}}_0^X}"'] \arrow[r,"f"]
\arrow[ur,phantom,"{\substack{\text{cocart et}\\\text{cart.}}}" description]
\& Y \arrow[u,hook]
\end{tikzcd},}
\[
i_f=f'\circ\mathrm{inj}_1^X=f'i_Xi_1^X
=j\,\underset{\substack{\text{inclusion can.}\\\text{de }X\text{ dans}\\Y\amalg X}}{\underline{(f\amalg\mathrm{id}_X)i_1^X}}
\]
est imm. ouverte de Dwyer, comme composée de $X\longrightarrow Y\amalg X$ et $j:Y\amalg X\longrightarrow Z(f)$ qui le sont ($j$ \inserted{obtenu} par cochangement de base à partir de $i_X$, qui est imm. ouverte de Dwyer).

\deleted{$\unsure{X\times I}$} D'autre part $X\xrightarrow{\mathrm{inj}_0^X}X\times I$ est imm. ouverte de Dwyer et dans $W$, donc couniversellement dans $W$, donc $Y\longrightarrow Z(f)$ est elle aussi dans $W$, donc aussi la rétraction $p_f$.

D'autre part, la proposition duale de
\[
\underline{\mathrm{Fib}}_W S\cap W_S=W^{\unsure{u|v}}_{S\,\unsure{\mathrm{univ.}|\mathrm{uni.}}}
\]
est valable pour les $W$-cofibrations :
% archive sheet number: 280
